\section{为什么学习深度强化学习}

\subsection{超越监督学习}

许多真实世界的 AI 应用不只是简单的输入输出映射。讲者给出几个关键场景：

\begin{enumerate}
    \item \textbf{决策有后果}：如 Spotify 推荐歌曲后，下一次推荐应该避免重复同一首歌，且应保持风格连续性。
    \item \textbf{直接监督不可得}：如编码助手（Cursor、Copilot），用户只能给出"好"或"不好"的反馈，而非直接提供正确输出。
    \item \textbf{目标不可微}：很多实际目标（如用户满意度）无法通过梯度直接优化。
\end{enumerate}

\begin{figure}[H]
\centering
\includegraphics[width=0.9\textwidth]{slides-images/slide-18.jpg}
\caption{课程用“部署会改变未来观测”这一视角解释为什么序贯决策问题不能退化为普通监督学习（Slide 18）。}
\end{figure}

\subsection{广泛应用于高性能 AI 系统}

讲者展示了大量深度 RL 的实际应用：

\begin{itemize}
    \item \textbf{机器人}：Unitree 的四足机器人行走、Physical Intelligence 的 $\pi_0$ 操作、Google Gemini Robotics。
    \item \textbf{游戏}：AlphaGo 的 Move 37 展示了 RL 发现人类未曾想到的新策略的能力。
    \item \textbf{大语言模型}：几乎所有现代 LLM 都使用某种形式的 RL 进行 post-training，尤其是在高级推理方面。
    \item \textbf{其他}：交通控制、生成式图像模型的 prompt 遵循、芯片设计（Google TPU）。
\end{itemize}

\begin{figure}[H]
\centering
\includegraphics[width=0.88\textwidth]{slides-images/slide-19.jpg}
\caption{深度 RL 已经成为复杂机器人系统的重要训练范式，尤其适合需要闭环控制和在线适应的物理任务（Slide 19）。}
\end{figure}

\begin{importantbox}{RL 的发现能力}
强化学习不只是模仿数据中已有的行为，它还能通过"试错"发现全新的解决方案。AlphaGo 的 Move 37 就是一个经典案例——它走出了人类棋手从未下过的棋步。
\end{importantbox}

\begin{figure}[H]
\centering
\includegraphics[width=0.88\textwidth]{slides-images/slide-22.jpg}
\caption{课程也明确把 LLM post-training 放进深度 RL 的应用版图中，说明 RL 已经从机器人和游戏扩展到语言模型系统（Slide 22）。}
\end{figure}

\subsection{学习的本质}

从哲学角度看，从经验中学习似乎是智能的基本组成部分。能够自主实践并改进的算法，是构建真正智能系统的关键。

\subsection{本章小结}

学习深度 RL 的理由包括：(1) 很多问题超出了监督学习的范畴；(2) 高性能 AI 系统广泛使用 RL；(3) 从经验中学习是智能的基本能力；(4) 领域中仍有大量开放研究问题。

